\documentclass{article}
\usepackage{iclr2027_conference,times}
%%%%% NEW MATH DEFINITIONS %%%%%

\usepackage{amsmath,amsfonts,bm}

% Mark sections of captions for referring to divisions of figures
\newcommand{\figleft}{{\em (Left)}}
\newcommand{\figcenter}{{\em (Center)}}
\newcommand{\figright}{{\em (Right)}}
\newcommand{\figtop}{{\em (Top)}}
\newcommand{\figbottom}{{\em (Bottom)}}
\newcommand{\captiona}{{\em (a)}}
\newcommand{\captionb}{{\em (b)}}
\newcommand{\captionc}{{\em (c)}}
\newcommand{\captiond}{{\em (d)}}

% Highlight a newly defined term
\newcommand{\newterm}[1]{{\bf #1}}


% Figure reference, lower-case.
\def\figref#1{figure~\ref{#1}}
% Figure reference, capital. For start of sentence
\def\Figref#1{Figure~\ref{#1}}
\def\twofigref#1#2{figures \ref{#1} and \ref{#2}}
\def\quadfigref#1#2#3#4{figures \ref{#1}, \ref{#2}, \ref{#3} and \ref{#4}}
% Section reference, lower-case.
\def\secref#1{section~\ref{#1}}
% Section reference, capital.
\def\Secref#1{Section~\ref{#1}}
% Reference to two sections.
\def\twosecrefs#1#2{sections \ref{#1} and \ref{#2}}
% Reference to three sections.
\def\secrefs#1#2#3{sections \ref{#1}, \ref{#2} and \ref{#3}}
% Reference to an equation, lower-case.
\def\eqref#1{equation~\ref{#1}}
% Reference to an equation, upper case
\def\Eqref#1{Equation~\ref{#1}}
% A raw reference to an equation---avoid using if possible
\def\plaineqref#1{\ref{#1}}
% Reference to a chapter, lower-case.
\def\chapref#1{chapter~\ref{#1}}
% Reference to an equation, upper case.
\def\Chapref#1{Chapter~\ref{#1}}
% Reference to a range of chapters
\def\rangechapref#1#2{chapters\ref{#1}--\ref{#2}}
% Reference to an algorithm, lower-case.
\def\algref#1{algorithm~\ref{#1}}
% Reference to an algorithm, upper case.
\def\Algref#1{Algorithm~\ref{#1}}
\def\twoalgref#1#2{algorithms \ref{#1} and \ref{#2}}
\def\Twoalgref#1#2{Algorithms \ref{#1} and \ref{#2}}
% Reference to a part, lower case
\def\partref#1{part~\ref{#1}}
% Reference to a part, upper case
\def\Partref#1{Part~\ref{#1}}
\def\twopartref#1#2{parts \ref{#1} and \ref{#2}}

\def\ceil#1{\lceil #1 \rceil}
\def\floor#1{\lfloor #1 \rfloor}
\def\1{\bm{1}}
\newcommand{\train}{\mathcal{D}}
\newcommand{\valid}{\mathcal{D_{\mathrm{valid}}}}
\newcommand{\test}{\mathcal{D_{\mathrm{test}}}}

\def\eps{{\epsilon}}


% Random variables
\def\reta{{\textnormal{$\eta$}}}
\def\ra{{\textnormal{a}}}
\def\rb{{\textnormal{b}}}
\def\rc{{\textnormal{c}}}
\def\rd{{\textnormal{d}}}
\def\re{{\textnormal{e}}}
\def\rf{{\textnormal{f}}}
\def\rg{{\textnormal{g}}}
\def\rh{{\textnormal{h}}}
\def\ri{{\textnormal{i}}}
\def\rj{{\textnormal{j}}}
\def\rk{{\textnormal{k}}}
\def\rl{{\textnormal{l}}}
% rm is already a command, just don't name any random variables m
\def\rn{{\textnormal{n}}}
\def\ro{{\textnormal{o}}}
\def\rp{{\textnormal{p}}}
\def\rq{{\textnormal{q}}}
\def\rr{{\textnormal{r}}}
\def\rs{{\textnormal{s}}}
\def\rt{{\textnormal{t}}}
\def\ru{{\textnormal{u}}}
\def\rv{{\textnormal{v}}}
\def\rw{{\textnormal{w}}}
\def\rx{{\textnormal{x}}}
\def\ry{{\textnormal{y}}}
\def\rz{{\textnormal{z}}}

% Random vectors
\def\rvepsilon{{\mathbf{\epsilon}}}
\def\rvtheta{{\mathbf{\theta}}}
\def\rva{{\mathbf{a}}}
\def\rvb{{\mathbf{b}}}
\def\rvc{{\mathbf{c}}}
\def\rvd{{\mathbf{d}}}
\def\rve{{\mathbf{e}}}
\def\rvf{{\mathbf{f}}}
\def\rvg{{\mathbf{g}}}
\def\rvh{{\mathbf{h}}}
\def\rvu{{\mathbf{i}}}
\def\rvj{{\mathbf{j}}}
\def\rvk{{\mathbf{k}}}
\def\rvl{{\mathbf{l}}}
\def\rvm{{\mathbf{m}}}
\def\rvn{{\mathbf{n}}}
\def\rvo{{\mathbf{o}}}
\def\rvp{{\mathbf{p}}}
\def\rvq{{\mathbf{q}}}
\def\rvr{{\mathbf{r}}}
\def\rvs{{\mathbf{s}}}
\def\rvt{{\mathbf{t}}}
\def\rvu{{\mathbf{u}}}
\def\rvv{{\mathbf{v}}}
\def\rvw{{\mathbf{w}}}
\def\rvx{{\mathbf{x}}}
\def\rvy{{\mathbf{y}}}
\def\rvz{{\mathbf{z}}}

% Elements of random vectors
\def\erva{{\textnormal{a}}}
\def\ervb{{\textnormal{b}}}
\def\ervc{{\textnormal{c}}}
\def\ervd{{\textnormal{d}}}
\def\erve{{\textnormal{e}}}
\def\ervf{{\textnormal{f}}}
\def\ervg{{\textnormal{g}}}
\def\ervh{{\textnormal{h}}}
\def\ervi{{\textnormal{i}}}
\def\ervj{{\textnormal{j}}}
\def\ervk{{\textnormal{k}}}
\def\ervl{{\textnormal{l}}}
\def\ervm{{\textnormal{m}}}
\def\ervn{{\textnormal{n}}}
\def\ervo{{\textnormal{o}}}
\def\ervp{{\textnormal{p}}}
\def\ervq{{\textnormal{q}}}
\def\ervr{{\textnormal{r}}}
\def\ervs{{\textnormal{s}}}
\def\ervt{{\textnormal{t}}}
\def\ervu{{\textnormal{u}}}
\def\ervv{{\textnormal{v}}}
\def\ervw{{\textnormal{w}}}
\def\ervx{{\textnormal{x}}}
\def\ervy{{\textnormal{y}}}
\def\ervz{{\textnormal{z}}}

% Random matrices
\def\rmA{{\mathbf{A}}}
\def\rmB{{\mathbf{B}}}
\def\rmC{{\mathbf{C}}}
\def\rmD{{\mathbf{D}}}
\def\rmE{{\mathbf{E}}}
\def\rmF{{\mathbf{F}}}
\def\rmG{{\mathbf{G}}}
\def\rmH{{\mathbf{H}}}
\def\rmI{{\mathbf{I}}}
\def\rmJ{{\mathbf{J}}}
\def\rmK{{\mathbf{K}}}
\def\rmL{{\mathbf{L}}}
\def\rmM{{\mathbf{M}}}
\def\rmN{{\mathbf{N}}}
\def\rmO{{\mathbf{O}}}
\def\rmP{{\mathbf{P}}}
\def\rmQ{{\mathbf{Q}}}
\def\rmR{{\mathbf{R}}}
\def\rmS{{\mathbf{S}}}
\def\rmT{{\mathbf{T}}}
\def\rmU{{\mathbf{U}}}
\def\rmV{{\mathbf{V}}}
\def\rmW{{\mathbf{W}}}
\def\rmX{{\mathbf{X}}}
\def\rmY{{\mathbf{Y}}}
\def\rmZ{{\mathbf{Z}}}

% Elements of random matrices
\def\ermA{{\textnormal{A}}}
\def\ermB{{\textnormal{B}}}
\def\ermC{{\textnormal{C}}}
\def\ermD{{\textnormal{D}}}
\def\ermE{{\textnormal{E}}}
\def\ermF{{\textnormal{F}}}
\def\ermG{{\textnormal{G}}}
\def\ermH{{\textnormal{H}}}
\def\ermI{{\textnormal{I}}}
\def\ermJ{{\textnormal{J}}}
\def\ermK{{\textnormal{K}}}
\def\ermL{{\textnormal{L}}}
\def\ermM{{\textnormal{M}}}
\def\ermN{{\textnormal{N}}}
\def\ermO{{\textnormal{O}}}
\def\ermP{{\textnormal{P}}}
\def\ermQ{{\textnormal{Q}}}
\def\ermR{{\textnormal{R}}}
\def\ermS{{\textnormal{S}}}
\def\ermT{{\textnormal{T}}}
\def\ermU{{\textnormal{U}}}
\def\ermV{{\textnormal{V}}}
\def\ermW{{\textnormal{W}}}
\def\ermX{{\textnormal{X}}}
\def\ermY{{\textnormal{Y}}}
\def\ermZ{{\textnormal{Z}}}

% Vectors
\def\vzero{{\bm{0}}}
\def\vone{{\bm{1}}}
\def\vmu{{\bm{\mu}}}
\def\vtheta{{\bm{\theta}}}
\def\va{{\bm{a}}}
\def\vb{{\bm{b}}}
\def\vc{{\bm{c}}}
\def\vd{{\bm{d}}}
\def\ve{{\bm{e}}}
\def\vf{{\bm{f}}}
\def\vg{{\bm{g}}}
\def\vh{{\bm{h}}}
\def\vi{{\bm{i}}}
\def\vj{{\bm{j}}}
\def\vk{{\bm{k}}}
\def\vl{{\bm{l}}}
\def\vm{{\bm{m}}}
\def\vn{{\bm{n}}}
\def\vo{{\bm{o}}}
\def\vp{{\bm{p}}}
\def\vq{{\bm{q}}}
\def\vr{{\bm{r}}}
\def\vs{{\bm{s}}}
\def\vt{{\bm{t}}}
\def\vu{{\bm{u}}}
\def\vv{{\bm{v}}}
\def\vw{{\bm{w}}}
\def\vx{{\bm{x}}}
\def\vy{{\bm{y}}}
\def\vz{{\bm{z}}}

% Elements of vectors
\def\evalpha{{\alpha}}
\def\evbeta{{\beta}}
\def\evepsilon{{\epsilon}}
\def\evlambda{{\lambda}}
\def\evomega{{\omega}}
\def\evmu{{\mu}}
\def\evpsi{{\psi}}
\def\evsigma{{\sigma}}
\def\evtheta{{\theta}}
\def\eva{{a}}
\def\evb{{b}}
\def\evc{{c}}
\def\evd{{d}}
\def\eve{{e}}
\def\evf{{f}}
\def\evg{{g}}
\def\evh{{h}}
\def\evi{{i}}
\def\evj{{j}}
\def\evk{{k}}
\def\evl{{l}}
\def\evm{{m}}
\def\evn{{n}}
\def\evo{{o}}
\def\evp{{p}}
\def\evq{{q}}
\def\evr{{r}}
\def\evs{{s}}
\def\evt{{t}}
\def\evu{{u}}
\def\evv{{v}}
\def\evw{{w}}
\def\evx{{x}}
\def\evy{{y}}
\def\evz{{z}}

% Matrix
\def\mA{{\bm{A}}}
\def\mB{{\bm{B}}}
\def\mC{{\bm{C}}}
\def\mD{{\bm{D}}}
\def\mE{{\bm{E}}}
\def\mF{{\bm{F}}}
\def\mG{{\bm{G}}}
\def\mH{{\bm{H}}}
\def\mI{{\bm{I}}}
\def\mJ{{\bm{J}}}
\def\mK{{\bm{K}}}
\def\mL{{\bm{L}}}
\def\mM{{\bm{M}}}
\def\mN{{\bm{N}}}
\def\mO{{\bm{O}}}
\def\mP{{\bm{P}}}
\def\mQ{{\bm{Q}}}
\def\mR{{\bm{R}}}
\def\mS{{\bm{S}}}
\def\mT{{\bm{T}}}
\def\mU{{\bm{U}}}
\def\mV{{\bm{V}}}
\def\mW{{\bm{W}}}
\def\mX{{\bm{X}}}
\def\mY{{\bm{Y}}}
\def\mZ{{\bm{Z}}}
\def\mBeta{{\bm{\beta}}}
\def\mPhi{{\bm{\Phi}}}
\def\mLambda{{\bm{\Lambda}}}
\def\mSigma{{\bm{\Sigma}}}

% Tensor
\DeclareMathAlphabet{\mathsfit}{\encodingdefault}{\sfdefault}{m}{sl}
\SetMathAlphabet{\mathsfit}{bold}{\encodingdefault}{\sfdefault}{bx}{n}
\newcommand{\tens}[1]{\bm{\mathsfit{#1}}}
\def\tA{{\tens{A}}}
\def\tB{{\tens{B}}}
\def\tC{{\tens{C}}}
\def\tD{{\tens{D}}}
\def\tE{{\tens{E}}}
\def\tF{{\tens{F}}}
\def\tG{{\tens{G}}}
\def\tH{{\tens{H}}}
\def\tI{{\tens{I}}}
\def\tJ{{\tens{J}}}
\def\tK{{\tens{K}}}
\def\tL{{\tens{L}}}
\def\tM{{\tens{M}}}
\def\tN{{\tens{N}}}
\def\tO{{\tens{O}}}
\def\tP{{\tens{P}}}
\def\tQ{{\tens{Q}}}
\def\tR{{\tens{R}}}
\def\tS{{\tens{S}}}
\def\tT{{\tens{T}}}
\def\tU{{\tens{U}}}
\def\tV{{\tens{V}}}
\def\tW{{\tens{W}}}
\def\tX{{\tens{X}}}
\def\tY{{\tens{Y}}}
\def\tZ{{\tens{Z}}}


% Graph
\def\gA{{\mathcal{A}}}
\def\gB{{\mathcal{B}}}
\def\gC{{\mathcal{C}}}
\def\gD{{\mathcal{D}}}
\def\gE{{\mathcal{E}}}
\def\gF{{\mathcal{F}}}
\def\gG{{\mathcal{G}}}
\def\gH{{\mathcal{H}}}
\def\gI{{\mathcal{I}}}
\def\gJ{{\mathcal{J}}}
\def\gK{{\mathcal{K}}}
\def\gL{{\mathcal{L}}}
\def\gM{{\mathcal{M}}}
\def\gN{{\mathcal{N}}}
\def\gO{{\mathcal{O}}}
\def\gP{{\mathcal{P}}}
\def\gQ{{\mathcal{Q}}}
\def\gR{{\mathcal{R}}}
\def\gS{{\mathcal{S}}}
\def\gT{{\mathcal{T}}}
\def\gU{{\mathcal{U}}}
\def\gV{{\mathcal{V}}}
\def\gW{{\mathcal{W}}}
\def\gX{{\mathcal{X}}}
\def\gY{{\mathcal{Y}}}
\def\gZ{{\mathcal{Z}}}

% Sets
\def\sA{{\mathbb{A}}}
\def\sB{{\mathbb{B}}}
\def\sC{{\mathbb{C}}}
\def\sD{{\mathbb{D}}}
% Don't use a set called E, because this would be the same as our symbol
% for expectation.
\def\sF{{\mathbb{F}}}
\def\sG{{\mathbb{G}}}
\def\sH{{\mathbb{H}}}
\def\sI{{\mathbb{I}}}
\def\sJ{{\mathbb{J}}}
\def\sK{{\mathbb{K}}}
\def\sL{{\mathbb{L}}}
\def\sM{{\mathbb{M}}}
\def\sN{{\mathbb{N}}}
\def\sO{{\mathbb{O}}}
\def\sP{{\mathbb{P}}}
\def\sQ{{\mathbb{Q}}}
\def\sR{{\mathbb{R}}}
\def\sS{{\mathbb{S}}}
\def\sT{{\mathbb{T}}}
\def\sU{{\mathbb{U}}}
\def\sV{{\mathbb{V}}}
\def\sW{{\mathbb{W}}}
\def\sX{{\mathbb{X}}}
\def\sY{{\mathbb{Y}}}
\def\sZ{{\mathbb{Z}}}

% Entries of a matrix
\def\emLambda{{\Lambda}}
\def\emA{{A}}
\def\emB{{B}}
\def\emC{{C}}
\def\emD{{D}}
\def\emE{{E}}
\def\emF{{F}}
\def\emG{{G}}
\def\emH{{H}}
\def\emI{{I}}
\def\emJ{{J}}
\def\emK{{K}}
\def\emL{{L}}
\def\emM{{M}}
\def\emN{{N}}
\def\emO{{O}}
\def\emP{{P}}
\def\emQ{{Q}}
\def\emR{{R}}
\def\emS{{S}}
\def\emT{{T}}
\def\emU{{U}}
\def\emV{{V}}
\def\emW{{W}}
\def\emX{{X}}
\def\emY{{Y}}
\def\emZ{{Z}}
\def\emSigma{{\Sigma}}

% entries of a tensor
% Same font as tensor, without \bm wrapper
\newcommand{\etens}[1]{\mathsfit{#1}}
\def\etLambda{{\etens{\Lambda}}}
\def\etA{{\etens{A}}}
\def\etB{{\etens{B}}}
\def\etC{{\etens{C}}}
\def\etD{{\etens{D}}}
\def\etE{{\etens{E}}}
\def\etF{{\etens{F}}}
\def\etG{{\etens{G}}}
\def\etH{{\etens{H}}}
\def\etI{{\etens{I}}}
\def\etJ{{\etens{J}}}
\def\etK{{\etens{K}}}
\def\etL{{\etens{L}}}
\def\etM{{\etens{M}}}
\def\etN{{\etens{N}}}
\def\etO{{\etens{O}}}
\def\etP{{\etens{P}}}
\def\etQ{{\etens{Q}}}
\def\etR{{\etens{R}}}
\def\etS{{\etens{S}}}
\def\etT{{\etens{T}}}
\def\etU{{\etens{U}}}
\def\etV{{\etens{V}}}
\def\etW{{\etens{W}}}
\def\etX{{\etens{X}}}
\def\etY{{\etens{Y}}}
\def\etZ{{\etens{Z}}}

% The true underlying data generating distribution
\newcommand{\pdata}{p_{\rm{data}}}
% The empirical distribution defined by the training set
\newcommand{\ptrain}{\hat{p}_{\rm{data}}}
\newcommand{\Ptrain}{\hat{P}_{\rm{data}}}
% The model distribution
\newcommand{\pmodel}{p_{\rm{model}}}
\newcommand{\Pmodel}{P_{\rm{model}}}
\newcommand{\ptildemodel}{\tilde{p}_{\rm{model}}}
% Stochastic autoencoder distributions
\newcommand{\pencode}{p_{\rm{encoder}}}
\newcommand{\pdecode}{p_{\rm{decoder}}}
\newcommand{\precons}{p_{\rm{reconstruct}}}

\newcommand{\laplace}{\mathrm{Laplace}} % Laplace distribution

\newcommand{\E}{\mathbb{E}}
\newcommand{\Ls}{\mathcal{L}}
\newcommand{\R}{\mathbb{R}}
\newcommand{\emp}{\tilde{p}}
\newcommand{\lr}{\alpha}
\newcommand{\reg}{\lambda}
\newcommand{\rect}{\mathrm{rectifier}}
\newcommand{\softmax}{\mathrm{softmax}}
\newcommand{\sigmoid}{\sigma}
\newcommand{\softplus}{\zeta}
\newcommand{\KL}{D_{\mathrm{KL}}}
\newcommand{\Var}{\mathrm{Var}}
\newcommand{\standarderror}{\mathrm{SE}}
\newcommand{\Cov}{\mathrm{Cov}}
% Wolfram Mathworld says $L^2$ is for function spaces and $\ell^2$ is for vectors
% But then they seem to use $L^2$ for vectors throughout the site, and so does
% wikipedia.
\newcommand{\normlzero}{L^0}
\newcommand{\normlone}{L^1}
\newcommand{\normltwo}{L^2}
\newcommand{\normlp}{L^p}
\newcommand{\normmax}{L^\infty}

\newcommand{\parents}{Pa} % See usage in notation.tex. Chosen to match Daphne's book.

\DeclareMathOperator*{\argmax}{arg\,max}
\DeclareMathOperator*{\argmin}{arg\,min}

\DeclareMathOperator{\sign}{sign}
\DeclareMathOperator{\Tr}{Tr}
\let\ab\allowbreak

\usepackage{amsmath,amssymb,booktabs,graphicx,hyperref,microtype,xcolor,xspace,url,needspace}
\hypersetup{hidelinks}
\newcommand{\method}{HR-Font\xspace}
\newcommand{\tc}{target-character appearance completion\xspace}
\newcommand{\Bset}{\mathcal B}
\newcommand{\Rset}{\mathcal R}

\title{HR-Font: Completing Font Families across Scripts}
\author{Anonymous authors}

\begin{document}
\maketitle

\begin{abstract}
We introduce \method, a framework for extending a font family to new scripts from a few Chinese reference glyphs.
The central challenge is to translate a family's visual language into characters whose structures differ from those observed.
\method addresses this challenge through two complementary forms of conditioning.
The Character Realization Prior expresses same-character examples from a font bank as anchor-relative structural variations and learns to combine them as generation unfolds.
Target-character appearance completion reads multi-scale local reference features to predict a spatial appearance representation of the requested glyph, supervised by training-target features.
The resulting conditions guide a diffusion generator toward recognizable characters with a coherent family identity.
We also formulate a cross-script completion benchmark that evaluates character identity, reference consistency, and family compatibility, with script-aware data curation and controlled reference budgets.
% Insert the quantitative conclusion after matched generation results and the
% evaluator review are complete. Missing measurements remain blank.
\end{abstract}

\section{Introduction}
\label{sec:intro}

Extending a typeface across writing systems requires more than transferring its visual style.
The target glyph must remain recognizable, while its proportions, curvature, stroke organization, terminals, and local details should remain consistent with the observed references.
Within a single script, this problem benefits from structural regularities shared by its characters.
Across scripts, the same visual attributes may need to be expressed through very different geometric forms.
A few Chinese references can reveal stroke weight, contrast, contour treatment, and decorative tendencies, but they do not determine how these attributes should appear in a Latin, kana, or bopomofo glyph.
Cross-script typeface extension is therefore not simply a matter of transferring the ``style'' of the references to target glyphs; it requires deciding how these style cues should interact with the target character's own structure to form a coherent design.

Most few-shot font generators follow a content--style conditioning formulation.
They improve synthesis through localized style representations, fine-grained content--reference interaction, richer content modeling, and stronger generative backbones
\citep{park2021lffont,park2021mxfont,tang2022fsfont,wang2023cffont,yang2024fontdiffuser}.
Cross-script methods extend this formulation with multi-level attention, multilingual features, and patch-level style modeling
\citep{li2021ftransgan,zhang2022mfnet,zhao2024fcagan,zeng2026rpfont,yan2026daafont}.
These methods improve the transfer of reference appearance across scripts, while leaving the structural realization of the target glyph implicit.
The same reference appearance can correspond to different widths, proportions, curves, stroke layouts, and local geometries for different target characters.
A model may therefore reproduce the overall appearance of the references while assigning it to a target structure that is not well matched to the intended design.

We model this ambiguity through two complementary factors.
For a fixed character identity, existing typefaces exhibit a range of geometric realizations.
These variations define an empirical deformation space for the character.
The source references provide the visual evidence needed to determine which realization is appropriate and how its local details should be rendered.
This separates the problem into two parts: how the target geometry changes, and how the resulting geometry expresses the observed appearance.

We propose HR-Font, a diffusion-based model built around this decomposition.
\emph{Character Realization Prior} (CRP) models character-specific geometric variation relative to a fixed neutral realization and produces a style-conditioned structural offset.
The neutral glyph specifies character identity, while the offset captures font-dependent changes in geometry.
\emph{Target-Character Appearance Completion} (TC) predicts a target-aligned appearance representation from the target content and source-script references.
It models how local properties such as contour treatment, stroke endings, texture, and decoration should appear on the requested target character.
The two conditions are combined with content and global reference features in the diffusion generator.

Because the same reference style may admit multiple plausible realizations in the target script, similarity to a single ground-truth glyph does not fully capture cross-script consistency.
We therefore report conventional image-quality and character metrics together with an independently trained cross-script style evaluator that measures consistency between source-script references and generated target glyphs.

Our contributions are:
\begin{itemize}
    \item We formulate cross-script font generation as a structure--appearance alignment problem and evaluate both character fidelity and cross-script consistency.
    \item We introduce \emph{Character Realization Prior} (CRP), which models target-character geometry through a style-conditioned deformation relative to a neutral realization.
    \item We introduce \emph{Target-Character Appearance Completion} (TC), which predicts how source-script appearance is realized on the local structure of a target character.
\end{itemize}

\begin{figure}[t]
\centering
% Editable LaTeX schematic; labels and equations are the source of truth.
\begingroup\small
\setlength{\tabcolsep}{6pt}
\begin{tabular}{p{0.25\linewidth}p{0.62\linewidth}}
\toprule
\textbf{Inputs} & \textbf{Conditions for the generator}\\
\midrule
Neutral target glyph &
$x_c^0 \longrightarrow E_c(x_c^0)$ \quad Character identity\\[5pt]
Same-character bank \newline + reference similarity &
$\displaystyle \Delta_c^s=\sum_j\alpha_j
 [E_c^s(B_{j,c})-E_c^s(x_c^0)]$\newline
Mean-Delta: anchor-relative variation\\[6pt]
Reference appearance \newline + target-content query &
$\{\psi(r_i)\},\,u_c \longrightarrow H \longrightarrow \hat a_{f,c}$\newline
TC: predicted target-character appearance\\[6pt]
\midrule
\textbf{Synthesis} &
$G'_j=G_j+W_o\hat a_{f,c}$\newline
Content + Mean-Delta + completed global / local references\newline
$\longrightarrow$ diffusion denoising $\longrightarrow \hat y_{f,c}$\\[5pt]
\bottomrule
\end{tabular}
\endgroup

\caption{Two complementary conditions for cross-script completion. CRP selects and combines anchor-relative structural variations from a font bank. TC reads local reference evidence into a target-aligned appearance memory. A diffusion generator uses both conditions alongside neutral content and global reference features. Training-target features supervise the appearance memory.}
\label{fig:overview}
\end{figure}

\section{Related Work}
\label{sec:related}

\subsection{Few-Shot Font Generation}

Few-shot font generation (FFG) aims to synthesize a complete typeface from a small number of reference glyphs.
A common approach is to separate character content from typeface style and improve the representation of either factor.
LF-Font models style with localized component-wise representations, while MX-Font uses multiple localized experts to capture reusable local style patterns
\citep{park2021lffont,park2021mxfont}.
FSFont further models fine-grained spatial correspondence between target content and reference glyphs, allowing local reference features to be aggregated according to the target character
\citep{tang2022fsfont}.
These methods establish the importance of localized representations for preserving detailed font appearance.

Another line of work focuses on content representation and generative modeling.
CF-Font enriches content features using corresponding glyphs from multiple basis fonts
\citep{wang2023cffont}.
FontDiffuser formulates font synthesis as conditional diffusion and combines multi-scale content aggregation with style contrastive learning
\citep{yang2024fontdiffuser}.
Recent approaches further explore global-aware autoregressive generation and explicit glyph spatial modeling
\citep{cai2026garfont,su2026glyphspatialnet}.
These methods improve content fidelity, structural quality, and local detail in few-shot generation.
HR-Font considers the related but different setting in which the structural regularities of the reference script do not directly determine the realization of a target-script character.

\subsection{Cross-Script Font Generation}

Cross-script font generation transfers typeface characteristics between writing systems with different character structures.
FTransGAN explicitly studies few-shot transfer between languages and uses multi-level attention to capture global and local reference information
\citep{li2021ftransgan}.
MF-Net extends multilingual font generation to previously unseen styles and languages using separate content and style encoders together with language-aware structural connections
\citep{zhang2022mfnet}.
FCAGAN improves cross-lingual style extraction through full-domain convolutional attention and multi-level feature modeling
\citep{zhao2024fcagan},
while DAA-Font combines dual-branch encoding with multi-level axial attention to model style and structural information across scripts
\citep{yan2026daafont}.
RPFont further improves cross-lingual transfer through patch-level style contrastive learning and relative-position-aware modeling
\citep{zeng2026rpfont}.
These methods primarily strengthen the representation and transfer of reference appearance across writing systems.

HR-Font additionally models the realization of the target character itself.
For a fixed character identity, typefaces can differ substantially in proportion, curvature, stroke organization, and local geometry.
Our Character Realization Prior (CRP) represents these font-dependent geometric variations relative to a neutral realization and provides them as a target-specific structural condition.
This separates the variation of target geometry from the appearance information inferred from the source-script references.

\subsection{Font Evaluation and Style Representation}

Font generation is commonly evaluated using complementary measures of reconstruction quality, perceptual similarity, distributional quality, and character recognizability.
SSIM measures structural image similarity
\citep{wang2004ssim},
LPIPS evaluates perceptual similarity in learned feature spaces
\citep{zhang2018lpips},
and FID compares the distributions of real and generated images
\citep{heusel2017fid}.
Font-generation studies have also used content-aware and style-aware classifiers to separately assess character correctness and style preservation
\citep{park2021lffont}.
These measures characterize different aspects of generation quality, but most compare images within a common script or against a particular target realization.

Learned font-style representations provide another way to measure typeface similarity.
Hassan et al.\ learn a font-style embedding space with metric learning, placing glyphs with similar font styles closer in the representation space
\citep{hassan2024fontstylespace}.
Relatedly, RPFont uses patch-level contrastive learning to learn stylistic representations that transfer across languages
\citep{zeng2026rpfont}.
These works show that font appearance can be compared in a learned representation space rather than only through image-level similarity.
For cross-script evaluation, we train an independent evaluator that learns a shared typeface representation across scripts and measures the consistency between source-script references and generated target glyphs.

\section{Cross-Script Font Completion}
\label{sec:task}
Let $c$ be a requested target character and $x_c^0$ its rendering in a neutral font.
For a target family $f$, the observed reference set is
$\Rset_f=\{(a_i,r_i)\}_{i=1}^{n}$, where $a_i$ is a Chinese character and $r_i$ its rendering in $f$.
The task is to generate $\hat y_{f,c}$ with the identity of $c$ and an appearance compatible with $f$.
The evaluation font is held out from training.
A bank $\Bset$ contains training-font renderings of the reference and target characters.
During training, an episode excludes its own font from donor retrieval.

\paragraph{Three evaluation dimensions.}
\emph{Identity} measures whether the output represents the requested character.
\emph{Reference consistency} measures agreement with visible attributes such as weight, contrast, stroke endings, and texture.
\emph{Family compatibility} measures whether the generated glyph is a coherent extension of the family across scripts.
The latter includes character-specific decisions whose exact realization is not shown in the reference set.
Paired target-image measurements complement these dimensions with reconstruction diagnostics.

\paragraph{Script-aware supervision.}
A font file can contain several scripts with different design conventions, fallback glyphs, or inconsistent coverage.
We curate valid font--character pairs and eligible donors by script, preserving the font-level train/validation/test split.
A mismatch in design conventions calls for family-level inspection; automatic similarity alone does not determine validity.
The released protocol records pair validity, donor eligibility, and the reference episode used for each output.

\section{Method}
\label{sec:method}

\subsection{Overview}
\label{sec:method-overview}

Cross-script typeface extension requires a target-character realization whose structure and local details are consistent with the reference family.
Chinese references reveal stroke treatment and decoration, while the neutral target provides one possible geometry.
We construct two target-character conditions.
The \emph{Character Realization Prior} (CRP) uses reference similarity to select candidate training families, then routes feature differences between their target-character renderings and the neutral glyph.
Two parallel branches convert these differences into structural offsets.
\emph{Target-Character Appearance Completion} (TC) predicts target-aligned local features from the references and supplies appearance tokens to the denoiser.

We use a FontDiffuser-based denoiser~\citep{yang2024fontdiffuser}, retaining its multi-scale content aggregation (MCA) and global-style conditioning.
A frozen content encoder $E_c$ provides target-content features, and a frozen style encoder $E_s$ provides global reference tokens $G$ and descriptors for candidate selection.
TC uses a separate trainable local encoder.
CRP modifies skip features through spatial offsets, while TC conditions decoder cross-attention.
Figure~\ref{fig:overview} summarizes the model.


\subsection{Character Realization Prior}
\label{sec:crp}

\paragraph{Reference-guided candidates.}
We compare the observed references with renderings of the same Chinese characters in a fixed pool of training families.
An admissible family contains the requested target character and all observed reference characters.
The episode font and its explicit weight variants are excluded before selection.
Using the pooled descriptor $g$ from $E_s$, we compute
\begin{equation}
\begin{aligned}
s_j
&=
\frac{1}{n}
\sum_{i=1}^{n}
\cos\!\left(
    g(r_i),
    g(x_{f_j,a_i})
\right),\\
\alpha_j
&=
\frac{\exp(s_j/\tau)}
     {\sum_{k\in\mathcal J}\exp(s_k/\tau)},
\qquad j\in\mathcal J.
\end{aligned}
\label{eq:prior}
\end{equation}
Here $\mathcal J=\mathcal J(c,\mathcal R_f)$ contains up to $K$ highest-scoring admissible families.
Each selected family provides a rendering $x_{f_j,c}$ of the requested target character.
The candidate set and prior $\alpha$ remain fixed throughout denoising.

\paragraph{Parallel feature differences.}
The two branches represent each candidate relative to the neutral target using spatial features from the frozen encoders:
\begin{equation}
\begin{aligned}
d^{\mathrm{ec},s}_{j,c}
&=
E_c^s(x_{f_j,c})-E_c^s(x_c^0),\\
d^{\mathrm{es},s}_{j,c}
&=
E_s^s(x_{f_j,c})-E_s^s(x_c^0).
\end{aligned}
\label{eq:delta}
\end{equation}
Here $s$ indexes the feature level.
At each structural module, the $E_s$ difference passes through a bias-free $1\times1$ adapter.
The branches share the candidates and prior weights, with separate routers and offset predictors.

\paragraph{Routing and structural offsets.}
Let $U_{\ell,t}$ be the skip feature at structural module $\ell$ and denoising step $t$.
Each branch forms a query $q^b_{\ell,t,p}$ at position $p$ from this feature, the neutral-content feature, the global reference summary, and the timestep.
Write $\bar d^b_{j,\ell}$ for the candidate difference at the matching feature level, including the adapter in the $E_s$ branch.
For $b\in\{\mathrm{ec},\mathrm{es}\}$,
\begin{equation}
\begin{aligned}
w^b_{j,\ell,t,p}
&=
\operatorname{softmax}_{j\in\mathcal J}\!\left[
    \log\alpha_j
    +\gamma^b_\ell
    \cos\!\left(
        q^b_{\ell,t,p},
        K^b_\ell(\bar d^b_{j,\ell})_p
    \right)
\right],\\
D^b_{\ell,t,p}
&=
\sum_{j\in\mathcal J}
w^b_{j,\ell,t,p}\bar d^b_{j,\ell,p},\\
\Delta^b_{\ell,t}
&=
F^b_\ell(U_{\ell,t},D^b_{\ell,t}),\\
\Delta_{\ell,t}
&=
\Delta^{\mathrm{ec}}_{\ell,t}
+\tanh(\zeta_\ell)\Delta^{\mathrm{es}}_{\ell,t}.
\end{aligned}
\label{eq:crp}
\end{equation}
$K^b_\ell$ is a key projection, $F^b_\ell$ is an offset predictor, and $\gamma^b_\ell$ and $\zeta_\ell$ are learned scalars.
Routing adapts the candidate mixture to each location and denoising state.
The branches predict offsets independently before the tanh-weighted sum.

The fused offset controls the residual skip update:
\begin{equation}
\widetilde U_{\ell,t}
=
U_{\ell,t}
+
P_\ell\!\left[
    \operatorname{DCN}_\ell(U_{\ell,t},\Delta_{\ell,t})
    -U_{\ell,t}
\right].
\label{eq:skip}
\end{equation}
$P_\ell$ is a zero-initialized $1\times1$ projection.
The updated skip is concatenated with the decoder feature before the corresponding residual block.


\subsection{Target-Character Appearance Completion}
\label{sec:tc}

TC predicts local appearance features at target-character positions from the neutral content and Chinese references.
A shared representation produces a spatial feature prediction supervised by the paired target and a separate token representation for the denoiser.

A shared trainable encoder $E_{\mathrm{loc}}$ extracts a shallow and a deeper feature map from each reference independently.
The shallow map is projected to the deeper map's channel width before both are flattened and concatenated into a token sequence $S_i$.
Target-content queries jointly read the reference sequences:
\begin{equation}
\begin{aligned}
\mathcal S
&=
[S_1;\ldots;S_n],
&
Q_c
&=
\mathcal Q(C_c)+\Pi_{\mathrm{slot}},\\
Z_0
&=
\operatorname{Attn}\!\left(
    Q_c,
    \mathcal K(\mathcal S),
    \mathcal V(\mathcal S)
\right),\\
Z
&=
Z_0+\operatorname{FFN}(\operatorname{LN}(Z_0)),
&
M
&=
\phi(Z),\quad T=\eta(Z).
\end{aligned}
\label{eq:tc}
\end{equation}
$C_c$ contains flattened target-content features from $E_c(x_c^0)$, $\Pi_{\mathrm{slot}}$ contains learned target-position vectors, and $\mathcal Q$, $\mathcal K$, and $\mathcal V$ are affine projections.
Attention uses head-wise query/key normalization.
For each query and head, its softmax spans all valid tokens across references, scales, and spatial positions.
Independent affine projections $\phi$ and $\eta$ produce the spatial prediction $M$ and denoiser tokens $T$.

We supervise $M$ with spatial VGG features~\citep{Simonyan15} of the paired target $y=x_{f,c}$:
\begin{equation}
\begin{aligned}
M^*
&=
\operatorname{sg}\!\left[
    \mathcal N\!\left(
        \operatorname{Pool}(V_2(y))
    \right)
\right],\\
\mathcal L_{\mathrm{comp}}
&=
\operatorname{SmoothL1}(M,M^*).
\end{aligned}
\label{eq:completion}
\end{equation}
$V_2$ denotes the frozen VGG feature stage used for local supervision.
Pooling matches the target-query grid and flattens it into rows, $\mathcal N$ applies fixed channel-wise normalization, and $\operatorname{sg}$ stops gradients through the target features.

\paragraph{Decoder conditioning.}
At each decoder cross-attention module, global tokens $G$ and local tokens $T$ are read with shared attention parameters and separate softmax normalizations.
For an active reference condition,
\begin{equation}
\begin{aligned}
A^G_{\ell,t}
&=
\operatorname{CA}_\ell(H_{\ell,t},G),
&
A^T_{\ell,t}
&=
\operatorname{CA}_\ell(H_{\ell,t},T),\\
\widetilde A_{\ell,t}
&=
A^G_{\ell,t}
+\beta_0r(u)\kappa_{\ell,t}A^T_{\ell,t},
\end{aligned}
\label{eq:inject}
\end{equation}
where $H_{\ell,t}$ is the current attention input.
The factor $\kappa_{\ell,t}$ is the detached, clipped global-to-local ratio of the outputs' root-mean-square (RMS) magnitudes.
$\beta_0$ is a fixed coefficient, and $r(u)$ ramps the local contribution over training updates $u$.


\subsection{Training and Inference}
\label{sec:training-inference}

\paragraph{Optimization.}
Starting from a pretrained denoiser, we jointly train it with both structural branches and the local reader.
The encoders $E_c$, $E_s$, and VGG remain frozen.
Each episode pairs one to eight Chinese references with a target glyph from the same family.
The objective is
\begin{equation}
\mathcal L
=
\mathcal L_{\mathrm{diff}}
+0.01\mathcal L_{\mathrm{perc}}
+r(u)\!\left(
    0.01\mathcal L_{\mathrm{comp}}
    +0.05\mathcal L_{\mathrm{detail}}
\right).
\label{eq:objective}
\end{equation}
$\mathcal L_{\mathrm{diff}}$ is the mean squared error of the predicted noise.
The perceptual loss compares VGG features of the paired target and the clean image reconstructed from the noise prediction.
The detail loss combines errors in target-defined regions, regions where the design adds or removes ink relative to the neutral rendering, and local high-frequency structure.
Architecture settings, calibration and ramp definitions, condition-dropout rules, and exact loss definitions appear in Appendix~\ref{app:k7b_impl}.

\paragraph{Inference.}
The candidate set, prior, and TC tokens are computed once per content-reference pair and reused throughout sampling.
CRP routing, decoder attention, and RMS calibration are recomputed at each denoising step.
The ramp is set to one at inference.
The denoiser predicts noise, and diffusion sampling produces the final glyph.

\section{Experiments}
\label{sec:experiments}
\subsection{Data and settings}
The corpus contains 260 fonts partitioned into 228 training, 16 validation, and 16 test fonts.
The inventory comprises 295 target characters and 338 Chinese reference characters.
Protocol A renders RGB glyphs at $96\times96$, with a shared font size within each font and neutral Noto Sans CJK Regular content.
The current clean training mapping retains 56,429 target pairs.
The stratified evaluation panel spans 47 characters covering Latin, diacritics, kana, and phonetic symbols; validity filtering retains 704 test pairs.

Reference-budget evaluation uses nested sets of 1, 2, 4, and 8 characters.
All methods in a comparison receive the same ordered reference set, and CRP retrieval uses that episode's references.
The main reference manifest and checkpoint selection are finalized on validation data.
The appendix records reference characters, bank eligibility, sampler, and provenance.

\paragraph{Training settings.}
The current joint experiments use eight V100 GPUs, eight samples per device, FP16 arithmetic, and gradient clipping at 1.0.
The effective main-task batch size is 64.
The denoiser and inherited structural modules use a peak learning rate of $2\times10^{-5}$; the local reader and router use $10^{-4}$.
Joint training runs for 20,000 successful updates, with 500 warmup updates, a plateau through update 5,000, and cosine decay to 10\% of the peak rate by update 10,000, held for the remaining updates.
The auxiliary experiment adds 32 Chinese examples per update and starts independently from the same parent weights.
EMA weights are evaluated every 2,000 updates.
Complete resolved settings appear in Appendix~\ref{app:repro}.

\subsection{Comparisons and measurements}
External comparisons cover conditional diffusion (FontDiffuser), cross-language transfer (FTransGAN), font-bank content fusion (CF-Font), and fine-grained reference attention (FSFont).
Each method is accompanied by its reference budget, bank access, pretraining data, and cross-script adaptation.
Internal comparisons isolate dynamic structural combination, local appearance completion, and auxiliary supervision using shared initial weights and data protocols.
The strongest validation baseline is retained in the main comparison.

We report character recognition accuracy, blind paired style preference, and family compatibility.
Human comparison displays the neutral target and the same reference set with randomized method positions and an explicit tie option.
The learned compatibility score uses a separately trained cross-script encoder and cosine similarity between a generated glyph and a Chinese-reference prototype.
We distinguish its font-retrieval validation from agreement with human judgments of generated appearance.
L1, SSIM, and LPIPS provide paired reconstruction diagnostics.
Results are aggregated by font and script over the same valid examples.

\begin{table}[t]
\centering\small
\caption{Main comparison under the cross-script protocol. Identity and style preference are percentages; family compatibility uses the cross-script cosine score. Empty cells await completed evaluation.}
\label{tab:main}
\begin{tabular}{lrrrr}
\toprule
Method & Identity $\uparrow$ & Style pref.\ $\uparrow$ & Family $\uparrow$ & LPIPS $\downarrow$\\
\midrule
FontDiffuser & -- & -- & -- & --\\
FTransGAN & -- & -- & -- & --\\
CF-Font & -- & -- & -- & --\\
FSFont & -- & -- & -- & --\\
CRP (fixed weights) & -- & -- & -- & --\\
\method & -- & -- & -- & --\\
\bottomrule
\end{tabular}
\end{table}

\subsection{Contribution studies}
\paragraph{Glyph variation prior.}
We compare routed combination with the retrieval prior used as fixed weights.
This comparison preserves bank access, anchor-relative features, and the appearance interface.
A same-bank absolute-feature control tests anchor-relative conditioning, while a shared-schedule structural-source comparison establishes the benefit of bank-supported variations.

\paragraph{Target-character appearance completion.}
We compare the full reader with a same-capacity branch trained using the generation objective alone.
A learned-slot-only query tests the role of target content, and a frozen-local-encoder variant tests adaptation of reference extraction.
These comparisons retain the 144-token interface, reference pixels, and training schedule.

\paragraph{Complementarity.}
Shared-parent variants separately add the structural and appearance components, followed by the full model.
The Chinese auxiliary study then measures whether additional same-family completion improves cross-script synthesis.
All variants are evaluated with identical reference episodes and diffusion noise.
Inference-time condition removal provides a separate sensitivity analysis of the trained model.

\begin{table}[t]
\centering\small
\caption{Contribution studies. The execution plan prioritizes the full model and auxiliary study, followed by targeted shared-budget comparisons.}
\label{tab:ablations}
\begin{tabular}{p{0.24\linewidth}p{0.56\linewidth}r}
\toprule
Question & Matched comparison & Result\\
\midrule
Structural source & Bank-supported CRP versus reference-derived structure & --\\
Residual interface & Anchor-relative versus absolute bank features & --\\
Routed combination & Routed versus fixed weights, same candidates & --\\
Appearance supervision & Spatial teacher versus generation loss only & --\\
Character conditioning & Content-plus-slot versus slot-only queries & --\\
Reference adaptation & Trainable versus frozen local encoder & --\\
Complementarity & Shared-parent structural / appearance / full model & --\\
Auxiliary supervision & Full model with versus without Chinese completion & --\\
\bottomrule
\end{tabular}
\end{table}

\subsection{Reference budgets and qualitative analysis}
Reference-budget curves evaluate the same checkpoints at 1, 2, 4, and 8 shots.
They measure how additional observations affect identity, appearance, and family coherence.
Qualitative comparisons include ordinary stroke styles, decorative fonts, and difficult cross-script realizations.
For each selected row, all methods share the target, references, sampler, and noise.
Figure~\ref{fig:qualitative} reserves the comparison layout for the completed outputs.

\begin{figure}[t]
\centering
% Blank cells are intentional; replace with actual matched glyph images.
\begingroup\small
\setlength{\tabcolsep}{7pt}
\begin{tabular}{lccccc}
\toprule
Case & References & Baseline & Mean-Delta & HR-Font & Target\\
\midrule
Stroke style & \rule{0pt}{22pt} & & & &\\
Decoration & \rule{0pt}{22pt} & & & &\\
Difficult completion & \rule{0pt}{22pt} & & & &\\
\bottomrule
\end{tabular}
\endgroup

\caption{Qualitative comparison layout. Rows cover stroke style, decorative treatment, and difficult cross-script completion. Blank cells will contain matched generated glyphs and the corresponding target design.}
\label{fig:qualitative}
\end{figure}

\section{Discussion}
The two conditions provide different forms of information.
CRP brings target-character examples into the generator's content coordinates and adapts their combination during synthesis.
TC learns a target-specific spatial appearance prediction from observed references.
Their combination gives the generator a structural prior and a supervised appearance destination.
The contribution studies measure how these roles translate into image quality.

The current representation operates on $96\times96$ raster glyphs.
Spatial appearance supervision connects local reference evidence to target locations.
The denoiser integrates this evidence into complete strokes, terminals, and decorative patterns.
Bank coverage also determines the available examples for rare target designs.
Extending the formulation to higher-resolution glyphs and complete typographic layout is a natural next step.

\section{Conclusion}
\method frames cross-script font generation as coherent family completion.
CRP provides an adaptive bank-supported variation prior relative to neutral content, and target-character appearance completion predicts how local reference observations apply to the requested glyph.
The accompanying evaluation separates identity, reference consistency, and family compatibility.
Together, the formulation and two conditioning components establish a practical framework for extending a font's visual language across scripts.
% Add the empirical conclusion after the result tables are populated.

\subsection*{Reproducibility statement}
The supplement records data mappings, cache fingerprints, reference episodes, parent checkpoints, resolved training configurations, and evaluation scripts.
The paper's editable figure sources accompany the draft.
Final results will include measured training and inference costs.

\subsection*{Ethics statement}
Font completion can reduce repetitive design work and improve script coverage.
Font licenses and source provenance must govern research use and redistribution.
The release will document permitted assets and the consent, compensation, and review procedures for human evaluation.

\subsection*{AI use statement}
AI tools assisted with research discussion, code review, drafting, and figure layout.
Experimental measurements come from the recorded training and evaluation runs.
The authors review the resulting text and artifacts and take responsibility for the submission.

\bibliography{references}
\bibliographystyle{iclr2027_conference}
\appendix

\section{Reproducibility Details}
\label{app:repro}
\begin{table}[ht]
\centering\small
\caption{Current implementation settings. Final run hashes and measured costs are filled from the selected artifacts.}
\begin{tabular}{ll}
\toprule
Setting & Value\\
\midrule
Image resolution & $96\times96$ RGB, Protocol A\\
Font split & 228 / 16 / 16\\
Valid training target pairs & 56,429\\
CRP retrieval & top-10, temperature 0.07\\
CRP router & width 32; location/module/time-conditioned weights\\
TC local evidence & $24^2+12^2=720$ tokens per reference\\
TC memory & 144 target queries, width 256, four attention heads\\
TC output / teacher & $144\times1024$ / $144\times128$\\
Reference injection & additive up-path attention; gain initialized at 0.01\\
Optimizer & AdamW, $(0.9,0.999)$, $\epsilon=10^{-8}$, weight decay 0.01\\
Joint batch / precision & $8\times8=64$ / FP16\\
Joint learning rates & inherited $2\times10^{-5}$; new $10^{-4}$\\
Updates & 20,000\\
Schedule & 500 warmup; plateau to 5k; cosine to 10\% at 10k, held to 20k\\
Gradient clipping & 1.0\\
Loss coefficients & perceptual 0.01; completion 0.01; detail 0.05; offset 0\\
Auxiliary Chinese study & batch 32; weight 0.5; shared main batch 64\\
Joint CFG / structural dropout & 0.02 / 0.05\\
Parent checkpoint hashes & --\\
Mapping and cache hashes & --\\
Training GPU hours / peak memory & --\\
Inference latency & --\\
\bottomrule
\end{tabular}
\end{table}

\paragraph{Implementation details.}
VGG inputs use ImageNet normalization; the teacher is adaptively pooled to $12\times12$ and standardized using clean training-target channel statistics.
Joint stages load the same parent checkpoint and create a new optimizer and scheduler.
The local encoder is a trainable copy of the first three global-style encoder blocks.
Reference attention uses normalized 64-dimensional query/key heads with FP32 logits and softmax.
CRP uses 32-dimensional routing features, with $1\times1$ hidden/content projections and a $3\times3$ candidate projection.
Its time embedding maps $[t/1000,(t/1000)^2]$ through a two-layer MLP; the routing strength starts at 0.1.
EMA decay rises to a maximum of 0.999, and the completion coefficient and local attention contribution ramp over 1,000 updates.
The structural offset is read at two levels of each encoder ($64\times48\times48$ and $128\times24\times24$) and is zero at the other levels.
The style-coordinate branch applies a bias-free $1\times1$ projection initialized to identity, and a per-module gate initialized at zero scales its contribution.
The detail term weights each sample by $\sqrt{\bar\alpha_t}$ and applies to samples whose condition is not dropped, and episodes draw one to eight references uniformly at random.

\paragraph{Run correspondence.}
The I0 reference is the 10,000-update G0b V0913 checkpoint.
I1 trains the full method; I2 adds Chinese auxiliary supervision and initializes independently from I0.
The Chinese auxiliary neutral inventory contains 338 glyphs rendered with the Protocol-A rule at font size 83; its images and frozen-content features have a separate manifest.
Chinese targets and references reuse the original training PNGs, and target characters are excluded from their reference sets.

\paragraph{Randomness and uncertainty.}
Training uses random state 3407.
Evaluation pairs reference episodes and diffusion noise across methods.
Planned 95\% intervals resample fonts while keeping paired method outputs together.
These intervals describe variation across evaluated fonts.

\paragraph{Reference episodes and sampler.}
The primary nested episode is specified by Unicode code points:
U+6C38, U+548C, U+4E66, U+98CE, U+9AA8, U+97F5, U+5929, U+5730.
The 1/2/4/8-shot sets use prefixes of this list.
The periodic validation panel uses 192 outputs at 1/4/8 shots.
The scheduled final validation panel comprises 16 validation fonts, 16 valid target characters per font, four reference episodes, and four budgets, yielding 4,096 outputs.
Inference uses DPM-Solver++ order 2, 20 steps, and guidance scale 7.5.
The larger 47-character protocol remains a separate evaluation panel with its own valid-pair manifest.

\paragraph{Model selection.}
Validation panels guide model and checkpoint selection using appearance, readability, and identity.
The final selection uses the full clean validation split.
Exploratory test visualizations are retained with their dates; subsequent selection uses validation data.
All reported methods are scored over an explicitly intersected valid sample set.

\section{Cross-Script Compatibility Evaluation}
\label{app:evaluator}
The current E12-b encoder $\phi$ is a ResNet-18 trained from random initialization on paired Chinese and Latin ground-truth glyphs.
It maps each $96\times96$ image to a unit-normalized 512-dimensional feature.
Training uses symmetric cross-script InfoNCE at temperature 0.07.
The scorer's data are drawn from 223 usable fonts in the generator's training partition; the generator's validation and test font files are excluded from scorer fitting.
The internal scorer split contains 156/33/34 font files for training, validation, and testing.
For a scoring reference set $\Rset_f$, define
\begin{equation}
 p_f=\frac{\sum_{r\in\Rset_f}\phi(r)}{\|\sum_{r\in\Rset_f}\phi(r)\|_2},
 \qquad S(\hat y,\Rset_f)=\frac{1+\langle\phi(\hat y),p_f\rangle}{2}.
\end{equation}
The score evaluates compatibility with the supplied references and requires no target image at scoring time.
Its retrieval study ranks Chinese reference prototypes using Latin glyph queries.
Its verification study compares same-font and different-font pairs.
The primary encoder study covers 52 Latin letters; script-specific validation accompanies extensions of this score to other target scripts.

A separate membership head combines query features with mean/max-pooled reference features and predicts same-font membership.
Its temperature is fitted on internal validation data; we treat this probability as a supplementary diagnostic.
Human evaluation measures visual agreement on generated outputs through blinded, paired judgments and font-clustered uncertainty.
Numerical evaluator validation and generation results will be added after the current protocol review and paired evaluation are finalized.

\clearpage
\section{Additional Results}
\begin{table}[!ht]
\centering\small
\caption{Reference-budget evaluation of the selected \method checkpoint.}
\begin{tabular}{lrrrr}
\toprule
References & Identity $\uparrow$ & Style pref.\ $\uparrow$ & Family $\uparrow$ & LPIPS $\downarrow$\\
\midrule
1 & -- & -- & -- & --\\
2 & -- & -- & -- & --\\
4 & -- & -- & -- & --\\
8 & -- & -- & -- & --\\
\bottomrule
\end{tabular}
\end{table}
% Add script-stratified tables, matched qualitative panels, and efficiency
% measurements from versioned evaluation outputs.
\end{document}
